\section{Acides aminés}

Les acides aminés ont une importance considérable, notamment dans le domaine de la biochimie. Outre des fonctions chimiques communes, leurs structures peuvent comporter des chaînes carbonées plus ou moins longues, éventuellement cycliques.

On s’intéresse à deux acides aminés linéaires à courtes chaînes carbonées : l’alanine et la glycine, dont les molécules sont représentées ci-dessous.


\begin{figure}[h]
    \centering
    \begin{minipage}[t]{0.45\textwidth}
        \centering
        \chemfig{OH-[:210,,1](=[:270]O)-[:150]-[:210,,,2]H_2N}
        \caption{Glycine}
    \end{minipage}%
    \hspace{0.05\textwidth}% Espace entre les deux figures
    \begin{minipage}[t]{0.45\textwidth}
        \centering
        \chemfig{OH-[:210,,1](=[:270]O)-[:150](-[:210,,,2]H_2N)(<[:75])(<:[:105]H)}
        \caption{Alanine}
    \end{minipage}
\end{figure}

\medskip
\noindent
\textbf{Q1.} Justifier l’appartenance de chacune de ces deux molécules à la famille des acides aminés en entourant et en nommant sur le \textbf{document réponse DR2 en page~\pageref{page:acide_amine_DR2}} (à rendre avec la copie), les groupes caractéristiques responsables.

\medskip
\noindent
\textbf{Q2.} L’une de ces deux molécules est chirale. L’identifier, avec justification.

\medskip
\noindent
\textbf{Q3.} Déterminer la configuration absolue correspondante en explicitant clairement votre démarche. On donne les numéros atomiques :
$Z(\ce{H}) = 1$, $Z(\ce{C}) = 6$, $Z(\ce{N}) = 7$ et $Z(\ce{O}) = 8$.

\medskip
\noindent
\textbf{Q4.} Représenter, en perspective de Cram, son énantiomère.

\medskip
La glycine, comme tous les acides aminés, est une espèce amphotère, c’est-à-dire acide et basique à la fois. Elle forme en solution aqueuse les espèces des deux couples acide-base suivants dont on fournit les $\mathrm{p}K_a$ à $\SI{20}{\celsius}$ :

\[
\begin{array}{l}
\mathrm{p}K_a \left(\ce{^+H_3N-CH_2-COOH} / \ce{^+H_3N-CH_2-COO^-}\right) = 2{,}4 \\
\mathrm{p}K_a \left(\ce{^+H_3N-CH_2-COO^-} / \ce{H_2N-CH_2-COO^-}\right) = 9{,}7 \\
\end{array}
\]

\medskip
\noindent
\textbf{Q5.} Écrire les deux équations modélisant l’échange de proton entre les deux espèces conjuguées de chacun des deux couples de la glycine.

\medskip
\noindent
\textbf{Q6.} Donner, pour chaque couple, l’expression de la relation liant les concentrations des espèces conjuguées du couple au $\mathrm{pH}$ de la solution et au $\mathrm{p}K_a$ du couple.

\medskip
\noindent
\textbf{Q7.} En déduire le diagramme de prédominance de la glycine. Vérifier qu’il comporte trois domaines correspondant aux trois espèces des deux couples acide-base auxquels appartient la glycine.

\medskip
\noindent
\textbf{Q8.} Déterminer, sans calculs, quelle est l’espèce prédominante de la glycine dans le sang à $\mathrm{pH} = 7{,}4$. En déduire une possible définition d’un zwitterion.

\medskip
\noindent
\textbf{Q9.} À $\SI{37}{\celsius}$, le produit ionique de l’eau dans le sang, noté $K_e$, est égal à $1{,}9 \times 10^{-14}$. En déduire le caractère acide ou basique du sang à cette température.

\cleardoublepage
\newpage
\newgeometry{top=1cm, bottom=1.5cm, left=1.5cm, right=1.5cm}
\begin{center}
\textbf{DOCUMENT RÉPONSE \\
À RENDRE OBLIGATOIREMENT AVEC LA COPIE}
\label{page:acide_amine_DR2}
\end{center}

\textbf{NOM :} \dotfill \textbf{Prénom :} \dotfill

\subsection*{Exercice – Acides aminés}

\medskip

\textbf{Document réponse DR2 :} Représentation des groupes caractéristiques de la glycine et de l’alanine.
\begin{figure}[h]
    \centering
    \begin{minipage}[t]{0.45\textwidth}
        \centering
        \chemfig{OH-[:210,,1](=[:270]O)-[:150]-[:210,,,2]H_2N}
        \caption{Glycine}
    \end{minipage}%
    \hspace{0.05\textwidth}% Espace entre les deux figures
    \begin{minipage}[t]{0.45\textwidth}
        \centering
        \chemfig{OH-[:210,,1](=[:270]O)-[:150](-[:210,,,2]H_2N)(<[:75])(<:[:105]H)}
        \caption{Alanine}
    \end{minipage}
\end{figure}

\restoregeometry
\newpage
